% Part: lambda-calculus
% Chapter: introduction
% Section: minimization

\documentclass[../../../include/open-logic-section]{subfiles}

\begin{document}

\olfileid{lam}{int}{min}
\olsection{توابع \usetoken{S}{lambda definable} تحت کمینه‌سازی بسته‌اند}

\begin{lem}
فرض کنید~$f(x,y)$ !!{lambda definable} باشد. $g$ را با
\[
g(x) \simeq \umin{y}{f(x,y)}.
\]
تعریف کنید. آنگاه~$g$ !!{lambda definable} است.
\end{lem}

\begin{proof}
اندیشهٔ کلی چنین است. با داشتن~$x$، از ترم نقطهٔ ثابت لامبدای~$Y$ برای
تعریف تابع~$h_x(n)$ استفاده می‌کنیم که جست‌وجوی~$y$ را از~$n$ آغاز
می‌کند؛ آن‌گاه~$g(x)$ همان~$h_x(0)$ است. تابع~$h_x$ را می‌توان به‌صورت
حل یک معادلهٔ نقطهٔ ثابت بیان کرد:
\[
h_x(n) \simeq
\begin{cases}
n & \text{اگر $f(x,n) = 0$} \\
h_x(n+1) & \text{در غیر این صورت.}
\end{cases}
\]

جزئیات چنین است. ازآنجاکه~$f$ تعریف‌پذیر با لامبداست، به‌وسیلهٔ ترمی
چون~$F$ !!{lambda defined} است. به یاد آورید که ترم لامبدای~$D$ را نیز
داریم، چنان‌که~$D(M, N, \bar{0}) \red M$ و~$D(M, N, \bar{1}) \red N$.
فعلاً~$x$ را ثابت بگیرید؛ برای اینکه ترمی~$h_x$ را !!{lambda define}،
می‌خواهیم ترم~$H$ای (وابسته به~$x$) بیابیم که
\[
H(\num{n}) \equiv D(\num{n}, H(S(\num{n})), F(x, \num{n})).
\]
را برآورده کند. این کار را می‌توان با ترم نقطهٔ ثابت~$Y$ انجام داد.
نخست~$U$ را ترم
\[
\lambd[h][\lambd[z][D(z,(h(Sz)),F(x,z))]],
\]
بگیرید و سپس~$H$ را ترم~$YU$ قرار دهید. توجه کنید که تنها متغیر آزاد
در~$H$، $x$ است. نشان می‌دهیم که~$H$ معادلهٔ بالا را برآورده می‌کند.

بنا بر تعریف~$Y$ داریم
\[
H = YU \equiv U(YU) = U(H).
\]
به‌ویژه، برای هر عدد طبیعی~$n$ داریم
\begin{align*}
H(\num{n}) & \equiv U(H, \num{n}) \\
& \red D(\num{n}, H(S(\num{n})), F(x, \num{n})),
\end{align*}
چنان‌که لازم بود. توجه کنید که اگر در سطر آخر عددنمای~$\num{m}$ را
به‌جای~$x$ بگذارید، عبارت هنگامی به~$\num{n}$ کاهش می‌یابد که
$F(\num{m}, \num{n})$ به~$\num{0}$ کاهش یابد؛ و اگر
$F(\num{m}, \num{n})$ به هر عددنمای دیگری کاهش یابد، عبارت به
$H(S(\num{n}))$ کاهش می‌یابد.

برای پایان‌دادن به اثبات، بگذارید~$G$ برابر~$\lambd[x][H(\num 0)]$ باشد.
آن‌گاه~$G$، تابع~$g$ را !!{lambda define}s؛ به بیان دیگر، برای هر~$m$،
$G(\num m)$ اگر~$g(m)$ تعریف‌شده باشد به~$\overline {g(m)}$ کاهش می‌یابد
و در غیر این صورت صورت نرمال ندارد.
\end{proof}

\end{document}
